\documentclass[11pt,a4paper]{article}
% lesson-preamble.tex -- Shared preamble for all Phase 00 lessons
% Usage: % lesson-preamble.tex -- Shared preamble for all Phase 00 lessons
% Usage: % lesson-preamble.tex -- Shared preamble for all Phase 00 lessons
% Usage: \input{../lesson-preamble} (from subdomain directories)

% ---- Encoding and fonts ----
\usepackage[utf8]{inputenc}
\usepackage[T1]{fontenc}

% ---- Mathematics ----
\usepackage{amsmath,amssymb,amsthm}
\usepackage{mathtools}

% ---- Page geometry ----
\usepackage[margin=1in,headheight=14pt]{geometry}

% ---- Hyperlinks ----
\usepackage[colorlinks=true,linkcolor=red,citecolor=blue,urlcolor=blue]{hyperref}

% ---- Lists ----
\usepackage{enumitem}

% ---- Colors and boxes ----
\usepackage{xcolor}
\usepackage[theorems,skins,breakable]{tcolorbox}

% ---- Header/Footer ----
\usepackage{fancyhdr}
\renewcommand{\headrulewidth}{0.4pt}

% ---- Tables ----
\usepackage{booktabs}

% ---- Bibliography ----
\usepackage{natbib}
\bibliographystyle{plainnat}

% --------------------------------------------------------------------------
% Theorem environments
% --------------------------------------------------------------------------
\theoremstyle{plain}
\newtheorem{theorem}{Theorem}[section]
\newtheorem{lemma}[theorem]{Lemma}
\newtheorem{proposition}[theorem]{Proposition}
\newtheorem{corollary}[theorem]{Corollary}

\theoremstyle{definition}
\newtheorem{definition}[theorem]{Definition}
\newtheorem{example}[theorem]{Example}
\newtheorem{exercise}[theorem]{Exercise}
\newtheorem{assumption}[theorem]{Assumption}

\theoremstyle{remark}
\newtheorem{remark}[theorem]{Remark}

% --------------------------------------------------------------------------
% Colored boxes
% --------------------------------------------------------------------------
\newtcolorbox{keyresult}[1][]{%
  colback=green!5!white,
  colframe=green!50!black,
  fonttitle=\bfseries,
  title={Key Result},
  breakable,
  #1
}

\newtcolorbox{rlconnection}[1][]{%
  colback=blue!5!white,
  colframe=blue!50!black,
  fonttitle=\bfseries,
  title={RL/ML Connection},
  breakable,
  #1
}

\newtcolorbox{warning}[1][]{%
  colback=red!5!white,
  colframe=red!50!black,
  fonttitle=\bfseries,
  title={Warning},
  breakable,
  #1
}

% --------------------------------------------------------------------------
% Custom commands -- number sets
% --------------------------------------------------------------------------
\newcommand{\R}{\mathbb{R}}
\newcommand{\N}{\mathbb{N}}
\newcommand{\Q}{\mathbb{Q}}
\newcommand{\Z}{\mathbb{Z}}
\newcommand{\C}{\mathbb{C}}
\newcommand{\F}{\mathbb{F}}

% --------------------------------------------------------------------------
% Custom commands -- probability and expectation
% --------------------------------------------------------------------------
\newcommand{\E}{\mathbb{E}}
\newcommand{\Prob}{\mathbb{P}}
\newcommand{\Var}{\operatorname{Var}}
\newcommand{\Cov}{\operatorname{Cov}}
\newcommand{\indicator}{\mathbf{1}}

% --------------------------------------------------------------------------
% Custom commands -- convergence arrows
% --------------------------------------------------------------------------
\newcommand{\cas}{\xrightarrow{\text{a.s.}}}
\newcommand{\cprob}{\xrightarrow{\Prob}}
\newcommand{\clp}[1]{\xrightarrow{L^{#1}}}
\newcommand{\cdist}{\xrightarrow{d}}

% --------------------------------------------------------------------------
% Custom commands -- common operators
% --------------------------------------------------------------------------
\DeclareMathOperator{\dom}{dom}
\DeclareMathOperator{\epi}{epi}
\DeclareMathOperator{\conv}{conv}
\DeclareMathOperator{\relint}{relint}
\DeclareMathOperator{\aff}{aff}
\DeclareMathOperator{\cl}{cl}
\DeclareMathOperator{\interior}{int}
\DeclareMathOperator{\tr}{tr}
\DeclareMathOperator{\diag}{diag}
\DeclareMathOperator{\rank}{rank}
\DeclareMathOperator{\sgn}{sgn}
\DeclareMathOperator{\supp}{supp}
\DeclareMathOperator{\argmin}{arg\,min}
\DeclareMathOperator{\argmax}{arg\,max}
\DeclareMathOperator{\prox}{prox}
\DeclareMathOperator{\dist}{dist}
\DeclareMathOperator{\diam}{diam}
\DeclareMathOperator{\Span}{span}

% --------------------------------------------------------------------------
% Custom commands -- linear algebra operators
% --------------------------------------------------------------------------
\DeclareMathOperator{\nullity}{nullity}
\DeclareMathOperator{\range}{range}
\DeclareMathOperator{\nullspace}{null}
\DeclareMathOperator{\proj}{proj}

% --------------------------------------------------------------------------
% Custom commands -- measure theory
% --------------------------------------------------------------------------
\newcommand{\powerset}{\mathcal{P}}
\newcommand{\Borel}{\mathcal{B}}
 (from subdomain directories)

% ---- Encoding and fonts ----
\usepackage[utf8]{inputenc}
\usepackage[T1]{fontenc}

% ---- Mathematics ----
\usepackage{amsmath,amssymb,amsthm}
\usepackage{mathtools}

% ---- Page geometry ----
\usepackage[margin=1in,headheight=14pt]{geometry}

% ---- Hyperlinks ----
\usepackage[colorlinks=true,linkcolor=red,citecolor=blue,urlcolor=blue]{hyperref}

% ---- Lists ----
\usepackage{enumitem}

% ---- Colors and boxes ----
\usepackage{xcolor}
\usepackage[theorems,skins,breakable]{tcolorbox}

% ---- Header/Footer ----
\usepackage{fancyhdr}
\renewcommand{\headrulewidth}{0.4pt}

% ---- Tables ----
\usepackage{booktabs}

% ---- Bibliography ----
\usepackage{natbib}
\bibliographystyle{plainnat}

% --------------------------------------------------------------------------
% Theorem environments
% --------------------------------------------------------------------------
\theoremstyle{plain}
\newtheorem{theorem}{Theorem}[section]
\newtheorem{lemma}[theorem]{Lemma}
\newtheorem{proposition}[theorem]{Proposition}
\newtheorem{corollary}[theorem]{Corollary}

\theoremstyle{definition}
\newtheorem{definition}[theorem]{Definition}
\newtheorem{example}[theorem]{Example}
\newtheorem{exercise}[theorem]{Exercise}
\newtheorem{assumption}[theorem]{Assumption}

\theoremstyle{remark}
\newtheorem{remark}[theorem]{Remark}

% --------------------------------------------------------------------------
% Colored boxes
% --------------------------------------------------------------------------
\newtcolorbox{keyresult}[1][]{%
  colback=green!5!white,
  colframe=green!50!black,
  fonttitle=\bfseries,
  title={Key Result},
  breakable,
  #1
}

\newtcolorbox{rlconnection}[1][]{%
  colback=blue!5!white,
  colframe=blue!50!black,
  fonttitle=\bfseries,
  title={RL/ML Connection},
  breakable,
  #1
}

\newtcolorbox{warning}[1][]{%
  colback=red!5!white,
  colframe=red!50!black,
  fonttitle=\bfseries,
  title={Warning},
  breakable,
  #1
}

% --------------------------------------------------------------------------
% Custom commands -- number sets
% --------------------------------------------------------------------------
\newcommand{\R}{\mathbb{R}}
\newcommand{\N}{\mathbb{N}}
\newcommand{\Q}{\mathbb{Q}}
\newcommand{\Z}{\mathbb{Z}}
\newcommand{\C}{\mathbb{C}}
\newcommand{\F}{\mathbb{F}}

% --------------------------------------------------------------------------
% Custom commands -- probability and expectation
% --------------------------------------------------------------------------
\newcommand{\E}{\mathbb{E}}
\newcommand{\Prob}{\mathbb{P}}
\newcommand{\Var}{\operatorname{Var}}
\newcommand{\Cov}{\operatorname{Cov}}
\newcommand{\indicator}{\mathbf{1}}

% --------------------------------------------------------------------------
% Custom commands -- convergence arrows
% --------------------------------------------------------------------------
\newcommand{\cas}{\xrightarrow{\text{a.s.}}}
\newcommand{\cprob}{\xrightarrow{\Prob}}
\newcommand{\clp}[1]{\xrightarrow{L^{#1}}}
\newcommand{\cdist}{\xrightarrow{d}}

% --------------------------------------------------------------------------
% Custom commands -- common operators
% --------------------------------------------------------------------------
\DeclareMathOperator{\dom}{dom}
\DeclareMathOperator{\epi}{epi}
\DeclareMathOperator{\conv}{conv}
\DeclareMathOperator{\relint}{relint}
\DeclareMathOperator{\aff}{aff}
\DeclareMathOperator{\cl}{cl}
\DeclareMathOperator{\interior}{int}
\DeclareMathOperator{\tr}{tr}
\DeclareMathOperator{\diag}{diag}
\DeclareMathOperator{\rank}{rank}
\DeclareMathOperator{\sgn}{sgn}
\DeclareMathOperator{\supp}{supp}
\DeclareMathOperator{\argmin}{arg\,min}
\DeclareMathOperator{\argmax}{arg\,max}
\DeclareMathOperator{\prox}{prox}
\DeclareMathOperator{\dist}{dist}
\DeclareMathOperator{\diam}{diam}
\DeclareMathOperator{\Span}{span}

% --------------------------------------------------------------------------
% Custom commands -- linear algebra operators
% --------------------------------------------------------------------------
\DeclareMathOperator{\nullity}{nullity}
\DeclareMathOperator{\range}{range}
\DeclareMathOperator{\nullspace}{null}
\DeclareMathOperator{\proj}{proj}

% --------------------------------------------------------------------------
% Custom commands -- measure theory
% --------------------------------------------------------------------------
\newcommand{\powerset}{\mathcal{P}}
\newcommand{\Borel}{\mathcal{B}}
 (from subdomain directories)

% ---- Encoding and fonts ----
\usepackage[utf8]{inputenc}
\usepackage[T1]{fontenc}

% ---- Mathematics ----
\usepackage{amsmath,amssymb,amsthm}
\usepackage{mathtools}

% ---- Page geometry ----
\usepackage[margin=1in,headheight=14pt]{geometry}

% ---- Hyperlinks ----
\usepackage[colorlinks=true,linkcolor=red,citecolor=blue,urlcolor=blue]{hyperref}

% ---- Lists ----
\usepackage{enumitem}

% ---- Colors and boxes ----
\usepackage{xcolor}
\usepackage[theorems,skins,breakable]{tcolorbox}

% ---- Header/Footer ----
\usepackage{fancyhdr}
\renewcommand{\headrulewidth}{0.4pt}

% ---- Tables ----
\usepackage{booktabs}

% ---- Bibliography ----
\usepackage{natbib}
\bibliographystyle{plainnat}

% --------------------------------------------------------------------------
% Theorem environments
% --------------------------------------------------------------------------
\theoremstyle{plain}
\newtheorem{theorem}{Theorem}[section]
\newtheorem{lemma}[theorem]{Lemma}
\newtheorem{proposition}[theorem]{Proposition}
\newtheorem{corollary}[theorem]{Corollary}

\theoremstyle{definition}
\newtheorem{definition}[theorem]{Definition}
\newtheorem{example}[theorem]{Example}
\newtheorem{exercise}[theorem]{Exercise}
\newtheorem{assumption}[theorem]{Assumption}

\theoremstyle{remark}
\newtheorem{remark}[theorem]{Remark}

% --------------------------------------------------------------------------
% Colored boxes
% --------------------------------------------------------------------------
\newtcolorbox{keyresult}[1][]{%
  colback=green!5!white,
  colframe=green!50!black,
  fonttitle=\bfseries,
  title={Key Result},
  breakable,
  #1
}

\newtcolorbox{rlconnection}[1][]{%
  colback=blue!5!white,
  colframe=blue!50!black,
  fonttitle=\bfseries,
  title={RL/ML Connection},
  breakable,
  #1
}

\newtcolorbox{warning}[1][]{%
  colback=red!5!white,
  colframe=red!50!black,
  fonttitle=\bfseries,
  title={Warning},
  breakable,
  #1
}

% --------------------------------------------------------------------------
% Custom commands -- number sets
% --------------------------------------------------------------------------
\newcommand{\R}{\mathbb{R}}
\newcommand{\N}{\mathbb{N}}
\newcommand{\Q}{\mathbb{Q}}
\newcommand{\Z}{\mathbb{Z}}
\newcommand{\C}{\mathbb{C}}
\newcommand{\F}{\mathbb{F}}

% --------------------------------------------------------------------------
% Custom commands -- probability and expectation
% --------------------------------------------------------------------------
\newcommand{\E}{\mathbb{E}}
\newcommand{\Prob}{\mathbb{P}}
\newcommand{\Var}{\operatorname{Var}}
\newcommand{\Cov}{\operatorname{Cov}}
\newcommand{\indicator}{\mathbf{1}}

% --------------------------------------------------------------------------
% Custom commands -- convergence arrows
% --------------------------------------------------------------------------
\newcommand{\cas}{\xrightarrow{\text{a.s.}}}
\newcommand{\cprob}{\xrightarrow{\Prob}}
\newcommand{\clp}[1]{\xrightarrow{L^{#1}}}
\newcommand{\cdist}{\xrightarrow{d}}

% --------------------------------------------------------------------------
% Custom commands -- common operators
% --------------------------------------------------------------------------
\DeclareMathOperator{\dom}{dom}
\DeclareMathOperator{\epi}{epi}
\DeclareMathOperator{\conv}{conv}
\DeclareMathOperator{\relint}{relint}
\DeclareMathOperator{\aff}{aff}
\DeclareMathOperator{\cl}{cl}
\DeclareMathOperator{\interior}{int}
\DeclareMathOperator{\tr}{tr}
\DeclareMathOperator{\diag}{diag}
\DeclareMathOperator{\rank}{rank}
\DeclareMathOperator{\sgn}{sgn}
\DeclareMathOperator{\supp}{supp}
\DeclareMathOperator{\argmin}{arg\,min}
\DeclareMathOperator{\argmax}{arg\,max}
\DeclareMathOperator{\prox}{prox}
\DeclareMathOperator{\dist}{dist}
\DeclareMathOperator{\diam}{diam}
\DeclareMathOperator{\Span}{span}

% --------------------------------------------------------------------------
% Custom commands -- linear algebra operators
% --------------------------------------------------------------------------
\DeclareMathOperator{\nullity}{nullity}
\DeclareMathOperator{\range}{range}
\DeclareMathOperator{\nullspace}{null}
\DeclareMathOperator{\proj}{proj}

% --------------------------------------------------------------------------
% Custom commands -- measure theory
% --------------------------------------------------------------------------
\newcommand{\powerset}{\mathcal{P}}
\newcommand{\Borel}{\mathcal{B}}

\pagestyle{fancy}
\fancyhf{}
\fancyhead[L]{\textsc{Lesson 10: Random Variables and Distributions}}
\fancyhead[R]{\thepage}
\title{Lesson 10: Random Variables and Distributions}
\author{AI/ML Mathematics}
\date{}
\begin{document}
\maketitle
\tableofcontents
\newpage

%% ============================================================
%% SECTION 1: MOTIVATION AND OVERVIEW
%% ============================================================

\section{Motivation and Overview}
\label{sec:motivation}

A reinforcement learning agent interacting with an environment receives
numerical rewards at each time step. The reward after pulling arm $k$ of a
multi-armed bandit might be $+1$ (success) or $0$ (failure); the total
number of successes in $n$ pulls could be any integer from $0$ to $n$; the
time until the first success is a positive integer; and the noise added to
a policy gradient estimate is often modeled as a continuous bell-shaped
quantity. Each of these quantities is a function that assigns a real number
to every possible experimental outcome. How should we formalize such
numerical summaries of random experiments? What tools let us compute the
probability that a random quantity falls in a given range, and what
standard families of distributions arise repeatedly in practice?

The tools developed in this lesson will let us formalize random variables
as measurable functions from a sample space to the real line, describe
their distributions through probability mass functions, probability
density functions, and cumulative distribution functions, and catalog the
discrete and continuous distributions most important for reinforcement
learning and machine learning.

\begin{rlconnection}[title={Why Random Variables and Distributions
    Matter for RL/ML}]
\begin{itemize}[nosep]
    \item \textbf{Bandit rewards}: Each arm of a multi-armed bandit
    produces Bernoulli-distributed rewards; the agent must estimate the
    mean of each arm's distribution.
    \item \textbf{Event counting in RL}: The number of times a rare event
    occurs in $n$ episodes is approximately Poisson-distributed, relevant
    for safety monitoring and anomaly detection.
    \item \textbf{Policy gradient noise}: The REINFORCE estimator adds
    Gaussian noise to gradient estimates; understanding the normal
    distribution is essential for variance analysis and confidence
    intervals.
\end{itemize}
\end{rlconnection}

\paragraph{Prerequisites.}
Probability spaces, Kolmogorov axioms, conditional probability, Bayes'
theorem, and independence (Lesson~9). Familiarity with sequences, series,
and basic integral calculus (Lessons~2--4).

\paragraph{Learning objectives.}
After completing this lesson, the student will be able to:
\begin{itemize}
    \item \emph{Define} random variables, probability mass functions,
    probability density functions, and cumulative distribution functions.
    \item \emph{Derive} the PMF, expectation, and variance of the
    Bernoulli, Binomial, Geometric, and Poisson distributions.
    \item \emph{Prove} the Poisson limit theorem (Poisson approximation to
    the Binomial).
    \item \emph{Compute} the PDF, CDF, expectation, and variance of the
    Uniform, Exponential, and Normal distributions.
    \item \emph{Prove} the Gaussian integral formula using polar
    coordinates.
    \item \emph{Explain} why $\Prob(X = x) = 0$ for continuous random
    variables and why density values are not probabilities.
\end{itemize}

%% ============================================================
%% SECTION 2: REFERENCES
%% ============================================================

\section{References}
\label{sec:references}

\begin{itemize}
    \item \citet[Chapters 3--5]{ross2019} -- Random variables, discrete
    and continuous distributions with many worked examples.
    \item \citet[Chapters 3--5]{blitzstein2019} -- Random variables and
    their distributions with a storytelling approach and extensive
    exercises.
    \item \citet[Chapters 2--4]{feller1968} -- Classical treatment of
    discrete distributions and generating functions.
    \item \citet[Chapter 2]{grimmett2020} -- Random variables, distribution
    functions, and density functions in a rigorous framework.
    \item \citet[Chapter 2]{bertsekas2008} -- Discrete and continuous
    random variables with engineering applications.
\end{itemize}

%% ============================================================
%% SECTION 3: RANDOM VARIABLES AND DISTRIBUTION FUNCTIONS
%% ============================================================

\section{Random Variables and Distribution Functions}
\label{sec:rv_cdf}

We begin by formalizing the notion of a numerical quantity determined by
the outcome of a random experiment.

\begin{definition}[Random Variable]
\label{def:random_variable}
Let $(\Omega, \Prob)$ be a probability space. A \emph{random variable} is
a function $X \colon \Omega \to \R$ that assigns a real number $X(\omega)$
to each outcome $\omega \in \Omega$.
\end{definition}

\begin{remark}
\label{rem:measurability}
In this lesson we work with discrete or ``well-behaved'' continuous sample
spaces where every function $X \colon \Omega \to \R$ is automatically
measurable. For uncountable sample spaces equipped with a
$\sigma$-algebra $\mathcal{F}$, one must additionally require that
$\{\omega : X(\omega) \leq x\} \in \mathcal{F}$ for every $x \in \R$;
this refinement is treated in Phase~01 (Lesson~1, measurable functions).
\end{remark}

\begin{definition}[Cumulative Distribution Function]
\label{def:cdf}
The \emph{cumulative distribution function} (CDF) of a random variable
$X$ is the function $F_X \colon \R \to [0,1]$ defined by
\[
    F_X(x) = \Prob(X \leq x) \quad \text{for all } x \in \R.
\]
\end{definition}

\begin{proposition}[Properties of the CDF]
\label{prop:cdf_properties}
For any random variable $X$, the CDF $F_X$ satisfies:
\begin{enumerate}[nosep,label=(\roman*)]
    \item $F_X$ is non-decreasing: if $x \leq y$, then
    $F_X(x) \leq F_X(y)$.
    \item $\lim_{x \to -\infty} F_X(x) = 0$ and
    $\lim_{x \to +\infty} F_X(x) = 1$.
    \item $F_X$ is right-continuous: $\lim_{h \to 0^+} F_X(x+h) = F_X(x)$
    for all $x \in \R$.
\end{enumerate}
\end{proposition}

\begin{proof}
\textbf{(i)} If $x \leq y$, then
$\{X \leq x\} \subseteq \{X \leq y\}$, so by monotonicity of $\Prob$,
$F_X(x) = \Prob(X \leq x) \leq \Prob(X \leq y) = F_X(y)$.

\textbf{(ii)} The events $A_n = \{X \leq -n\}$ satisfy
$A_1 \supseteq A_2 \supseteq \cdots$ and
$\bigcap_{n=1}^\infty A_n = \emptyset$. By continuity from above
(Lesson~9, Theorem~3.5),
$\lim_{n \to \infty} F_X(-n) = \Prob(\emptyset) = 0$.
Since $F_X$ is non-decreasing, $\lim_{x \to -\infty} F_X(x) = 0$.
Similarly, $B_n = \{X \leq n\}$ gives an increasing sequence with
$\bigcup B_n = \Omega$, so
$\lim_{x \to +\infty} F_X(x) = 1$ by continuity from below.

\textbf{(iii)} Let $h_n \downarrow 0$. Then
$\{X \leq x + h_n\} \downarrow \{X \leq x\}$ (since any $\omega$ with
$X(\omega) \leq x$ is in every $\{X \leq x + h_n\}$, and any $\omega$
with $X(\omega) > x$ satisfies $X(\omega) > x + h_n$ for $n$ large
enough). By continuity from above,
$\lim_n F_X(x + h_n) = F_X(x)$.
\end{proof}

\begin{example}[CDF of a Coin Flip]
\label{ex:cdf_coin}
Toss a fair coin. Define $X = 1$ if heads, $X = 0$ if tails. Then
\[
    F_X(x) = \begin{cases}
        0 & \text{if } x < 0, \\
        1/2 & \text{if } 0 \leq x < 1, \\
        1 & \text{if } x \geq 1.
    \end{cases}
\]
The CDF has jumps of size $1/2$ at $x = 0$ and $x = 1$.
\end{example}

We now distinguish two fundamental types of random variables: discrete and
continuous.

%% ============================================================
%% SECTION 4: DISCRETE RANDOM VARIABLES AND PMF
%% ============================================================

\section{Discrete Random Variables}
\label{sec:discrete}

\begin{definition}[Discrete Random Variable and PMF]
\label{def:discrete_rv}
A random variable $X$ is \emph{discrete} if it takes values in a
countable set $\{x_1, x_2, \ldots\}$. Its \emph{probability mass
function} (PMF) is
\[
    p_X(x) = \Prob(X = x) \quad \text{for each } x \in \R.
\]
The PMF satisfies $p_X(x) \geq 0$ for all $x$, $p_X(x) = 0$ for
$x \notin \{x_1, x_2, \ldots\}$, and $\sum_i p_X(x_i) = 1$.
\end{definition}

\begin{definition}[Expectation of a Discrete Random Variable]
\label{def:discrete_expectation}
Let $X$ be a discrete random variable with PMF $p_X$. The
\emph{expectation} (or \emph{mean}) of $X$ is
\[
    \E[X] = \sum_{x} x \, p_X(x),
\]
provided the sum converges absolutely, i.e.,
$\sum_x |x| \, p_X(x) < \infty$.
\end{definition}

\begin{definition}[Variance of a Discrete Random Variable]
\label{def:discrete_variance}
The \emph{variance} of a discrete random variable $X$ with
$\E[X^2] < \infty$ is
\[
    \Var(X) = \E\bigl[(X - \E[X])^2\bigr]
    = \sum_x (x - \E[X])^2 \, p_X(x).
\]
The \emph{standard deviation} is $\sigma_X = \sqrt{\Var(X)}$.
\end{definition}

\begin{proposition}[Computational Formula for Variance]
\label{prop:var_formula}
For any random variable $X$ with finite second moment,
\[
    \Var(X) = \E[X^2] - (\E[X])^2.
\]
\end{proposition}

\begin{proof}
Let $\mu = \E[X]$. Then
\begin{align*}
    \Var(X) &= \E[(X - \mu)^2]
    = \E[X^2 - 2\mu X + \mu^2] \\
    &= \E[X^2] - 2\mu \E[X] + \mu^2
    = \E[X^2] - 2\mu^2 + \mu^2
    = \E[X^2] - \mu^2. \qedhere
\end{align*}
\end{proof}

We now catalog the most important discrete distributions.

\subsection{Bernoulli Distribution}

\begin{definition}[Bernoulli Distribution]
\label{def:bernoulli}
A random variable $X$ has a \emph{Bernoulli distribution} with parameter
$p \in [0,1]$, written $X \sim \mathrm{Bernoulli}(p)$, if $X$ takes
values in $\{0, 1\}$ with PMF
\[
    p_X(k) = p^k (1-p)^{1-k}, \quad k \in \{0,1\}.
\]
That is, $\Prob(X = 1) = p$ and $\Prob(X = 0) = 1 - p$.
\end{definition}

\begin{proposition}[Bernoulli Moments]
\label{prop:bernoulli_moments}
If $X \sim \mathrm{Bernoulli}(p)$, then $\E[X] = p$ and
$\Var(X) = p(1-p)$.
\end{proposition}

\begin{proof}
Since $X \in \{0,1\}$, we have $X^2 = X$, so
$\E[X] = 0 \cdot (1-p) + 1 \cdot p = p$ and
$\E[X^2] = \E[X] = p$. By
Proposition~\ref{prop:var_formula},
$\Var(X) = p - p^2 = p(1-p)$.
\end{proof}

\begin{rlconnection}[title={Bernoulli Rewards in Multi-Armed Bandits}]
In a $K$-armed Bernoulli bandit, pulling arm $k$ yields a reward
$R_k \sim \mathrm{Bernoulli}(p_k)$. The agent's goal is to maximize the
total expected reward $\sum_{t=1}^T \E[R_{A_t}]$ where $A_t$ is the arm
selected at time $t$. The regret analysis of algorithms like UCB and
Thompson sampling relies on the Bernoulli mean $p_k$ and variance
$p_k(1 - p_k)$ to quantify the difficulty of distinguishing near-optimal
arms.
\end{rlconnection}

\begin{example}[Bernoulli Variance]
\label{ex:bernoulli_var}
The variance $p(1-p)$ is maximized at $p = 1/2$ (most uncertain case)
and equals zero at $p = 0$ or $p = 1$ (deterministic outcomes). This is
why a fair coin flip is the ``hardest'' Bernoulli trial to predict: it
has the largest spread.
\end{example}

\subsection{Binomial Distribution}

\begin{definition}[Binomial Distribution]
\label{def:binomial}
A random variable $X$ has a \emph{Binomial distribution} with parameters
$n \in \N$ and $p \in [0,1]$, written $X \sim \mathrm{Bin}(n, p)$, if
\[
    p_X(k) = \binom{n}{k} p^k (1-p)^{n-k}, \quad k = 0, 1, \ldots, n.
\]
\end{definition}

\begin{proposition}[Binomial PMF is Valid]
\label{prop:binomial_valid}
The Binomial PMF sums to $1$.
\end{proposition}

\begin{proof}
By the binomial theorem,
\[
    \sum_{k=0}^n \binom{n}{k} p^k (1-p)^{n-k}
    = \bigl(p + (1-p)\bigr)^n = 1^n = 1. \qedhere
\]
\end{proof}

\begin{theorem}[Binomial Moments via Indicator Decomposition]
\label{thm:binomial_moments}
If $X \sim \mathrm{Bin}(n, p)$, then $\E[X] = np$ and
$\Var(X) = np(1-p)$.
\end{theorem}

\begin{proof}
Write $X = X_1 + X_2 + \cdots + X_n$ where each
$X_i \sim \mathrm{Bernoulli}(p)$ independently represents the outcome of
the $i$-th trial. Since expectation is additive (regardless of
independence; see Lesson~11),
\[
    \E[X] = \sum_{i=1}^n \E[X_i] = np.
\]
For the variance, independence gives
$\Var(X) = \sum_{i=1}^n \Var(X_i) = n \cdot p(1-p)$.
(The variance of a sum of independent random variables equals the sum of
variances; this is proved in Lesson~11, Theorem~4.3.)
\end{proof}

\begin{example}[Quality Control]
\label{ex:binomial_quality}
A production line has a $2\%$ defect rate. In a batch of $n = 100$ items,
the number of defectives $X \sim \mathrm{Bin}(100, 0.02)$ has
$\E[X] = 2$ and $\Var(X) = 100 \cdot 0.02 \cdot 0.98 = 1.96$, so
$\sigma_X \approx 1.40$. Most batches will have between 0 and 5 defects.
\end{example}

\subsection{Geometric Distribution}

\begin{definition}[Geometric Distribution]
\label{def:geometric}
A random variable $X$ has a \emph{Geometric distribution} with parameter
$p \in (0, 1]$, written $X \sim \mathrm{Geom}(p)$, if $X$ counts the
number of the trial on which the first success occurs in a sequence of
independent Bernoulli$(p)$ trials. Its PMF is
\[
    p_X(k) = (1-p)^{k-1} p, \quad k = 1, 2, 3, \ldots
\]
\end{definition}

\begin{proposition}[Geometric PMF is Valid]
\label{prop:geometric_valid}
The Geometric PMF sums to $1$.
\end{proposition}

\begin{proof}
Let $q = 1 - p$. Then
\[
    \sum_{k=1}^\infty q^{k-1} p
    = p \sum_{j=0}^\infty q^j
    = p \cdot \frac{1}{1 - q}
    = \frac{p}{p} = 1,
\]
where the geometric series converges since $0 \leq q < 1$.
\end{proof}

\begin{proposition}[Geometric Moments]
\label{prop:geometric_moments}
If $X \sim \mathrm{Geom}(p)$, then $\E[X] = 1/p$ and
$\Var(X) = (1 - p)/p^2$.
\end{proposition}

\begin{proof}
Let $q = 1 - p$. For the expectation, we differentiate the geometric
series. Recall $\sum_{k=0}^\infty q^k = 1/(1-q)$. Differentiating both
sides with respect to $q$:
$\sum_{k=1}^\infty k q^{k-1} = 1/(1-q)^2 = 1/p^2$.
Therefore
\[
    \E[X] = \sum_{k=1}^\infty k \, q^{k-1} p
    = p \cdot \frac{1}{p^2} = \frac{1}{p}.
\]

For the second moment, differentiate $\sum_{k=1}^\infty k q^{k-1} =
1/p^2$ with respect to $q$:
$\sum_{k=2}^\infty k(k-1) q^{k-2} = 2/p^3$.
Then
\begin{align*}
    \E[X(X-1)] &= \sum_{k=1}^\infty k(k-1) q^{k-1} p
    = p \cdot q \cdot \frac{2}{p^3} = \frac{2q}{p^2}.
\end{align*}
So $\E[X^2] = \E[X(X-1)] + \E[X] = 2q/p^2 + 1/p = (2q + p)/p^2
= (q + 1)/p^2$. Using Proposition~\ref{prop:var_formula},
\[
    \Var(X) = \frac{q + 1}{p^2} - \frac{1}{p^2}
    = \frac{q}{p^2} = \frac{1-p}{p^2}. \qedhere
\]
\end{proof}

\begin{example}[Expected Waiting Time]
\label{ex:geometric_waiting}
A server processes requests, each succeeding independently with
probability $p = 0.3$. The expected number of attempts until the first
success is $\E[X] = 1/0.3 \approx 3.33$, with standard deviation
$\sigma_X = \sqrt{0.7/0.09} \approx 2.79$.
\end{example}

\subsection{Poisson Distribution}

\begin{definition}[Poisson Distribution]
\label{def:poisson}
A random variable $X$ has a \emph{Poisson distribution} with parameter
$\lambda > 0$, written $X \sim \mathrm{Poisson}(\lambda)$, if
\[
    p_X(k) = \frac{\lambda^k e^{-\lambda}}{k!}, \quad k = 0, 1, 2, \ldots
\]
\end{definition}

\begin{proposition}[Poisson PMF is Valid]
\label{prop:poisson_valid}
The Poisson PMF sums to $1$.
\end{proposition}

\begin{proof}
Using the Taylor series $e^\lambda = \sum_{k=0}^\infty \lambda^k / k!$,
\[
    \sum_{k=0}^\infty \frac{\lambda^k e^{-\lambda}}{k!}
    = e^{-\lambda} \sum_{k=0}^\infty \frac{\lambda^k}{k!}
    = e^{-\lambda} \cdot e^\lambda = 1. \qedhere
\]
\end{proof}

\begin{proposition}[Poisson Moments]
\label{prop:poisson_moments}
If $X \sim \mathrm{Poisson}(\lambda)$, then $\E[X] = \lambda$ and
$\Var(X) = \lambda$.
\end{proposition}

\begin{proof}
For the expectation:
\[
    \E[X] = \sum_{k=0}^\infty k \frac{\lambda^k e^{-\lambda}}{k!}
    = \sum_{k=1}^\infty \frac{\lambda^k e^{-\lambda}}{(k-1)!}
    = \lambda e^{-\lambda} \sum_{j=0}^\infty \frac{\lambda^j}{j!}
    = \lambda e^{-\lambda} \cdot e^\lambda = \lambda.
\]

For $\E[X(X-1)]$:
\[
    \E[X(X-1)] = \sum_{k=2}^\infty \frac{k(k-1) \lambda^k e^{-\lambda}}{k!}
    = \lambda^2 e^{-\lambda} \sum_{j=0}^\infty \frac{\lambda^j}{j!}
    = \lambda^2.
\]
Therefore $\E[X^2] = \lambda^2 + \lambda$, and
$\Var(X) = \E[X^2] - (\E[X])^2 = \lambda^2 + \lambda - \lambda^2
= \lambda$.
\end{proof}

\begin{keyresult}[title={Poisson Limit Theorem}]
The Poisson distribution arises as the limit of the Binomial distribution
when $n \to \infty$ and $p \to 0$ such that $np \to \lambda$. This
justifies using the Poisson as an approximation for counting rare events
in many independent trials.
\end{keyresult}

\begin{theorem}[Poisson Limit Theorem]
\label{thm:poisson_limit}
Let $X_n \sim \mathrm{Bin}(n, p_n)$ where $p_n = \lambda / n$ for a
fixed $\lambda > 0$. Then for every $k = 0, 1, 2, \ldots$,
\[
    \lim_{n \to \infty} \Prob(X_n = k)
    = \frac{\lambda^k e^{-\lambda}}{k!}.
\]
\end{theorem}

\begin{proof}
We have
\[
    \Prob(X_n = k) = \binom{n}{k} \left(\frac{\lambda}{n}\right)^k
    \left(1 - \frac{\lambda}{n}\right)^{n-k}.
\]
Write $\binom{n}{k} = n(n-1)\cdots(n-k+1)/k!$ and factor:
\begin{align*}
    \Prob(X_n = k)
    &= \frac{n(n-1)\cdots(n-k+1)}{k!} \cdot \frac{\lambda^k}{n^k}
       \cdot \left(1 - \frac{\lambda}{n}\right)^{n-k} \\
    &= \frac{\lambda^k}{k!} \cdot
       \underbrace{\frac{n}{n} \cdot \frac{n-1}{n}
       \cdots \frac{n-k+1}{n}}_{\to 1 \text{ as } n \to \infty}
       \cdot \underbrace{\left(1 - \frac{\lambda}{n}\right)^n}_{%
       \to e^{-\lambda}}
       \cdot \underbrace{\left(1 - \frac{\lambda}{n}\right)^{-k}}_{%
       \to 1}.
\end{align*}
The first underbraced factor is a product of $k$ terms each tending to
$1$. The second factor converges to $e^{-\lambda}$ by the well-known
limit $\lim_{n \to \infty} (1 - \lambda/n)^n = e^{-\lambda}$. The third
factor tends to $1$ since $k$ is fixed and
$(1 - \lambda/n) \to 1$. Combining:
$\lim_{n \to \infty} \Prob(X_n = k) = \lambda^k e^{-\lambda} / k!$.
\end{proof}

\begin{rlconnection}[title={Poisson Processes in RL Event Counting}]
In safety-critical RL applications, an agent may need to monitor the
occurrence of rare failure events across many episodes. If each episode
independently has a small failure probability $p$, the total failure
count across $n$ episodes is approximately $\mathrm{Poisson}(np)$ by the
Poisson limit theorem. This approximation simplifies the computation of
tail probabilities for alarm thresholds and risk-sensitive objectives.
\end{rlconnection}

\begin{example}[Poisson Approximation]
\label{ex:poisson_approx}
A website receives $n = 10{,}000$ visitors per hour, each clicking a
specific ad with probability $p = 0.0003$. The exact distribution of
clicks is $\mathrm{Bin}(10000, 0.0003)$, but the Poisson approximation
gives $\lambda = np = 3$. Then $\Prob(X = 0) \approx e^{-3} \approx
0.050$ and $\Prob(X \leq 2) \approx e^{-3}(1 + 3 + 9/2) \approx 0.423$.
\end{example}

Having covered the key discrete distributions, we now turn to continuous
random variables.

%% ============================================================
%% SECTION 5: CONTINUOUS RANDOM VARIABLES
%% ============================================================

\section{Continuous Random Variables}
\label{sec:continuous}

\begin{definition}[Continuous Random Variable and PDF]
\label{def:continuous_rv}
A random variable $X$ is \emph{continuous} if there exists a non-negative
function $f_X \colon \R \to [0, \infty)$, called the \emph{probability
density function} (PDF), such that for every interval $[a, b]$,
\[
    \Prob(a \leq X \leq b) = \int_a^b f_X(x) \, dx.
\]
The PDF satisfies $f_X(x) \geq 0$ for all $x$ and
$\int_{-\infty}^{\infty} f_X(x) \, dx = 1$.
\end{definition}

\begin{proposition}[CDF-PDF Relationship]
\label{prop:cdf_pdf}
If $X$ is a continuous random variable with PDF $f_X$, then
\[
    F_X(x) = \int_{-\infty}^x f_X(t) \, dt,
\]
and at every point where $f_X$ is continuous,
$f_X(x) = F_X'(x)$.
\end{proposition}

\begin{proof}
By definition,
$F_X(x) = \Prob(X \leq x) = \int_{-\infty}^x f_X(t) \, dt$.
By the Fundamental Theorem of Calculus, at points of continuity of $f_X$,
the derivative of this integral equals $f_X(x)$.
\end{proof}

\begin{warning}[title={$\Prob(X = x) = 0$ for Continuous Random Variables}]
For a continuous random variable, $\Prob(X = x) = 0$ for every
$x \in \R$. This is because
$\Prob(X = x) = \Prob(x \leq X \leq x) = \int_x^x f_X(t)\,dt = 0$.
In particular, $f_X(x)$ is \emph{not} the probability that $X$ equals
$x$; rather, $f_X(x)\,dx$ represents an infinitesimal probability.
Density values can exceed $1$ (e.g., the PDF of
$\mathrm{Uniform}(0, 1/2)$ equals $2$ on $(0, 1/2)$).
\end{warning}

\begin{definition}[Expectation and Variance of Continuous Random Variables]
\label{def:continuous_expectation}
If $X$ is continuous with PDF $f_X$, then
\[
    \E[X] = \int_{-\infty}^{\infty} x \, f_X(x) \, dx
\]
provided the integral converges absolutely, and
\[
    \Var(X) = \int_{-\infty}^{\infty} (x - \E[X])^2 \, f_X(x) \, dx
    = \E[X^2] - (\E[X])^2.
\]
\end{definition}

\subsection{Uniform Distribution}

\begin{definition}[Continuous Uniform Distribution]
\label{def:uniform}
A random variable $X$ has a \emph{Uniform distribution} on $[a, b]$ with
$a < b$, written $X \sim \mathrm{Uniform}(a, b)$, if its PDF is
\[
    f_X(x) = \begin{cases}
        \frac{1}{b - a} & \text{if } a \leq x \leq b, \\
        0 & \text{otherwise}.
    \end{cases}
\]
\end{definition}

\begin{proposition}[Uniform Moments]
\label{prop:uniform_moments}
If $X \sim \mathrm{Uniform}(a, b)$, then $\E[X] = (a + b)/2$ and
$\Var(X) = (b - a)^2/12$.
\end{proposition}

\begin{proof}
\[
    \E[X] = \int_a^b \frac{x}{b-a} \, dx
    = \frac{1}{b-a} \cdot \frac{x^2}{2}\bigg|_a^b
    = \frac{b^2 - a^2}{2(b-a)}
    = \frac{a + b}{2}.
\]
For the second moment,
\[
    \E[X^2] = \int_a^b \frac{x^2}{b-a} \, dx
    = \frac{1}{b-a} \cdot \frac{x^3}{3}\bigg|_a^b
    = \frac{b^3 - a^3}{3(b-a)}
    = \frac{a^2 + ab + b^2}{3}.
\]
Therefore
\begin{align*}
    \Var(X) &= \E[X^2] - (\E[X])^2
    = \frac{a^2 + ab + b^2}{3} - \frac{(a+b)^2}{4} \\
    &= \frac{4(a^2 + ab + b^2) - 3(a^2 + 2ab + b^2)}{12}
    = \frac{a^2 - 2ab + b^2}{12}
    = \frac{(b-a)^2}{12}. \qedhere
\end{align*}
\end{proof}

\begin{example}[Uniform Random Number Generator]
\label{ex:uniform_rng}
For $X \sim \mathrm{Uniform}(0, 1)$, we have $\E[X] = 1/2$ and
$\Var(X) = 1/12 \approx 0.083$. The CDF is $F_X(x) = x$ for
$x \in [0,1]$, which is the basis for the inverse transform method:
if $U \sim \mathrm{Uniform}(0,1)$ and $F$ is any CDF, then
$X = F^{-1}(U)$ has CDF $F$.
\end{example}

\subsection{Exponential Distribution}

\begin{definition}[Exponential Distribution]
\label{def:exponential}
A random variable $X$ has an \emph{Exponential distribution} with rate
parameter $\lambda > 0$, written $X \sim \mathrm{Exp}(\lambda)$, if its
PDF is
\[
    f_X(x) = \begin{cases}
        \lambda e^{-\lambda x} & \text{if } x \geq 0, \\
        0 & \text{if } x < 0.
    \end{cases}
\]
\end{definition}

\begin{proposition}[Exponential Moments]
\label{prop:exponential_moments}
If $X \sim \mathrm{Exp}(\lambda)$, then $\E[X] = 1/\lambda$ and
$\Var(X) = 1/\lambda^2$.
\end{proposition}

\begin{proof}
Using integration by parts with $u = x$, $dv = \lambda e^{-\lambda x}dx$:
\[
    \E[X] = \int_0^\infty x \lambda e^{-\lambda x} \, dx
    = \left[-x e^{-\lambda x}\right]_0^\infty
    + \int_0^\infty e^{-\lambda x} \, dx
    = 0 + \frac{1}{\lambda}
    = \frac{1}{\lambda}.
\]
For $\E[X^2]$, integrate by parts twice:
\[
    \E[X^2] = \int_0^\infty x^2 \lambda e^{-\lambda x} \, dx
    = \frac{2}{\lambda^2}.
\]
(Alternatively, $\E[X^2] = \E[X(X-1)] + \E[X]$; using the substitution
$j = k-2$ in the factorial moment computation gives the same result.)
Therefore $\Var(X) = 2/\lambda^2 - 1/\lambda^2 = 1/\lambda^2$.
\end{proof}

\begin{theorem}[Memoryless Property of the Exponential Distribution]
\label{thm:memoryless}
Let $X \sim \mathrm{Exp}(\lambda)$. Then for all $s, t \geq 0$,
\[
    \Prob(X > s + t \mid X > s) = \Prob(X > t).
\]
Moreover, the Exponential distribution is the unique continuous
distribution with the memoryless property.
\end{theorem}

\begin{proof}
The CDF is $F_X(x) = 1 - e^{-\lambda x}$ for $x \geq 0$, so
$\Prob(X > x) = e^{-\lambda x}$. By the definition of conditional
probability,
\[
    \Prob(X > s + t \mid X > s)
    = \frac{\Prob(X > s + t)}{\Prob(X > s)}
    = \frac{e^{-\lambda(s+t)}}{e^{-\lambda s}}
    = e^{-\lambda t}
    = \Prob(X > t).
\]

For uniqueness, suppose $X$ is a continuous non-negative random variable
satisfying $\Prob(X > s + t) = \Prob(X > s)\Prob(X > t)$ for all
$s, t \geq 0$. Let $g(t) = \Prob(X > t)$. Then $g(s + t) = g(s)g(t)$,
$g(0) = 1$, and $g$ is non-increasing and right-continuous (since
$F_X = 1 - g$ is a CDF). The only non-increasing solutions to
Cauchy's functional equation $g(s+t) = g(s)g(t)$ with $g(0) = 1$ that
are measurable (guaranteed by monotonicity) are $g(t) = e^{-\lambda t}$
for some $\lambda \geq 0$. Since $g(t) \to 0$ as $t \to \infty$
(because $X$ is a proper random variable), we need $\lambda > 0$.
This gives $X \sim \mathrm{Exp}(\lambda)$.
\end{proof}

\begin{example}[Exponential Waiting Time]
\label{ex:exponential_waiting}
A light bulb has an exponentially distributed lifetime with mean
$1/\lambda = 1000$ hours. Given that it has survived 500 hours, the
memoryless property says the conditional distribution of the remaining
lifetime is still $\mathrm{Exp}(\lambda)$ with mean 1000 hours. The
bulb's past survival provides no information about its future.
\end{example}

\subsection{Normal (Gaussian) Distribution}

\begin{definition}[Normal Distribution]
\label{def:normal}
A random variable $X$ has a \emph{Normal} (or \emph{Gaussian})
distribution with parameters $\mu \in \R$ and $\sigma^2 > 0$, written
$X \sim \mathcal{N}(\mu, \sigma^2)$, if its PDF is
\[
    f_X(x) = \frac{1}{\sqrt{2\pi\sigma^2}}
    \exp\!\left(-\frac{(x - \mu)^2}{2\sigma^2}\right),
    \quad x \in \R.
\]
The case $\mu = 0$, $\sigma^2 = 1$ is the \emph{standard normal}
distribution.
\end{definition}

To verify that this is a valid PDF, we must show that the integral over
$\R$ equals $1$. It suffices to evaluate the integral for the standard
normal.

\begin{keyresult}[title={The Gaussian Integral}]
The integral $\int_{-\infty}^{\infty} e^{-x^2/2}\,dx = \sqrt{2\pi}$
is the foundational identity that normalizes the Gaussian density. It
cannot be evaluated by elementary antiderivatives; the proof uses a
change to polar coordinates.
\end{keyresult}

\begin{theorem}[Gaussian Integral via Polar Coordinates]
\label{thm:gaussian_integral}
\[
    I = \int_{-\infty}^{\infty} e^{-x^2/2} \, dx = \sqrt{2\pi}.
\]
\end{theorem}

\begin{proof}
Since $e^{-x^2/2} > 0$, the integral $I$ is positive (and converges
because $e^{-x^2/2}$ decays faster than any polynomial). We compute
$I^2$ instead:
\[
    I^2 = \left(\int_{-\infty}^{\infty} e^{-x^2/2} \, dx\right)
    \left(\int_{-\infty}^{\infty} e^{-y^2/2} \, dy\right)
    = \int_{-\infty}^{\infty} \int_{-\infty}^{\infty}
    e^{-(x^2 + y^2)/2} \, dx \, dy.
\]
Switch to polar coordinates $x = r\cos\theta$, $y = r\sin\theta$ with
Jacobian $r$:
\begin{align*}
    I^2 &= \int_0^{2\pi} \int_0^{\infty} e^{-r^2/2} \, r \, dr \, d\theta
    = 2\pi \int_0^{\infty} r \, e^{-r^2/2} \, dr.
\end{align*}
Substitute $u = r^2/2$, $du = r\,dr$:
\[
    \int_0^{\infty} r \, e^{-r^2/2} \, dr
    = \int_0^{\infty} e^{-u} \, du
    = \left[-e^{-u}\right]_0^{\infty} = 1.
\]
Therefore $I^2 = 2\pi$, and since $I > 0$, we conclude
$I = \sqrt{2\pi}$.
\end{proof}

\begin{corollary}
\label{cor:normal_valid}
The standard normal PDF $f(x) = (2\pi)^{-1/2} e^{-x^2/2}$ integrates
to $1$ over $\R$. For general $\mathcal{N}(\mu, \sigma^2)$, the
substitution $z = (x - \mu)/\sigma$ with $dx = \sigma \, dz$ gives
\[
    \int_{-\infty}^{\infty} \frac{1}{\sqrt{2\pi\sigma^2}}
    e^{-(x-\mu)^2/(2\sigma^2)} dx
    = \int_{-\infty}^{\infty} \frac{1}{\sqrt{2\pi}} e^{-z^2/2} \, dz
    = 1.
\]
\end{corollary}

\begin{theorem}[Normal Moments via Standardization]
\label{thm:normal_moments}
If $X \sim \mathcal{N}(\mu, \sigma^2)$, then $\E[X] = \mu$ and
$\Var(X) = \sigma^2$.
\end{theorem}

\begin{proof}
\emph{Standard normal first.} Let $Z \sim \mathcal{N}(0,1)$. Since
$z \, e^{-z^2/2}$ is an odd function and $e^{-z^2/2}$ decays rapidly,
\[
    \E[Z] = \frac{1}{\sqrt{2\pi}} \int_{-\infty}^{\infty}
    z \, e^{-z^2/2} \, dz = 0.
\]
For the second moment, integrate by parts with $u = z$,
$dv = z \, e^{-z^2/2}\,dz$ (so $v = -e^{-z^2/2}$):
\begin{align*}
    \E[Z^2] &= \frac{1}{\sqrt{2\pi}} \int_{-\infty}^{\infty}
    z^2 e^{-z^2/2} \, dz
    = \frac{1}{\sqrt{2\pi}}
    \left(\left[-z \, e^{-z^2/2}\right]_{-\infty}^{\infty}
    + \int_{-\infty}^{\infty} e^{-z^2/2} \, dz\right) \\
    &= \frac{1}{\sqrt{2\pi}} \left(0 + \sqrt{2\pi}\right) = 1.
\end{align*}
Thus $\Var(Z) = \E[Z^2] - (\E[Z])^2 = 1 - 0 = 1$.

\emph{General case.} Write $X = \mu + \sigma Z$ where
$Z \sim \mathcal{N}(0,1)$. Then
$\E[X] = \mu + \sigma \E[Z] = \mu$ and
$\Var(X) = \sigma^2 \Var(Z) = \sigma^2$ (using the affine
transformation rule $\Var(aZ + b) = a^2 \Var(Z)$, proved in Lesson~11).
\end{proof}

\begin{rlconnection}[title={Gaussian Noise in Policy Gradient Methods}]
In continuous-action policy gradient methods such as REINFORCE, the
policy $\pi_\theta(a \mid s)$ is typically a Gaussian
$\mathcal{N}(\mu_\theta(s), \sigma_\theta^2(s))$. The gradient estimate
\[
    \nabla_\theta J(\theta)
    \approx \frac{1}{N} \sum_{i=1}^N
    \nabla_\theta \log \pi_\theta(a_i \mid s_i) \, G_i
\]
involves the log-density of the normal distribution. Understanding the
Gaussian PDF, its mean, and variance is essential for computing
$\nabla_\theta \log \pi_\theta$, analyzing estimator variance, and
designing variance reduction baselines.
\end{rlconnection}

\begin{example}[Standard Normal Tail Probabilities]
\label{ex:normal_tails}
For $Z \sim \mathcal{N}(0,1)$, the ``68-95-99.7 rule'' states
\[
    \Prob(|Z| \leq 1) \approx 0.683, \quad
    \Prob(|Z| \leq 2) \approx 0.954, \quad
    \Prob(|Z| \leq 3) \approx 0.997.
\]
This means observations more than $3\sigma$ from the mean are extremely
rare under Gaussian models, a fact used in outlier detection and
confidence interval construction.
\end{example}

%% ============================================================
%% SECTION 6: EXERCISES
%% ============================================================

\section{Exercises}
\label{sec:exercises}

\begin{exercise}[Binomial Probabilities]
\textnormal{(\(\ast\))} \\
A fair coin is tossed $n = 6$ times. Let $X$ be the number of heads.
\begin{enumerate}[nosep,label=(\alph*)]
    \item Compute $\Prob(X = 3)$.
    \item Compute $\Prob(X \geq 5)$.
    \item Compute $\E[X]$ and $\Var(X)$ using the Binomial formulas.
\end{enumerate}
\end{exercise}

\begin{exercise}[Geometric Distribution and the Memoryless Property]
\textnormal{(\(\ast\))} \\
Let $X \sim \mathrm{Geom}(p)$ count the trial of the first success.
\begin{enumerate}[nosep,label=(\alph*)]
    \item Prove that the Geometric distribution has the discrete
    memoryless property: $\Prob(X > m + n \mid X > m) = \Prob(X > n)$
    for all $m, n \geq 0$.
    \item Is the Geometric distribution the unique discrete distribution
    with this property? Prove or disprove.
\end{enumerate}
\end{exercise}

\begin{exercise}[Poisson Approximation Error]
\textnormal{(\(\ast\ast\))} \\
Let $X \sim \mathrm{Bin}(n, \lambda/n)$ and $Y \sim
\mathrm{Poisson}(\lambda)$ with $\lambda = 2$ and $n = 20$.
\begin{enumerate}[nosep,label=(\alph*)]
    \item Compute $\Prob(X = 0)$, $\Prob(X = 1)$, and $\Prob(X = 2)$
    exactly using the Binomial PMF.
    \item Compute the corresponding Poisson approximations
    $\Prob(Y = 0)$, $\Prob(Y = 1)$, $\Prob(Y = 2)$.
    \item Quantify the approximation error for each value.
\end{enumerate}
\end{exercise}

\begin{exercise}[Exponential CDF and Quantiles]
\textnormal{(\(\ast\ast\))} \\
Let $X \sim \mathrm{Exp}(\lambda)$.
\begin{enumerate}[nosep,label=(\alph*)]
    \item Derive the CDF $F_X(x)$ from the PDF.
    \item Find the median of $X$ (i.e., the value $m$ such that
    $F_X(m) = 1/2$).
    \item Show that the median is always less than the mean $1/\lambda$.
    What does this tell us about the shape of the exponential
    distribution?
\end{enumerate}
\end{exercise}

\begin{exercise}[Gaussian Moments via Differentiation]
\textnormal{(\(\ast\ast\ast\))} \\
Define $I(\alpha) = \int_{-\infty}^{\infty} e^{-\alpha x^2} \, dx$
for $\alpha > 0$.
\begin{enumerate}[nosep,label=(\alph*)]
    \item Using the Gaussian integral (Theorem~\ref{thm:gaussian_integral}
    with the substitution $u = \sqrt{2\alpha}\,x$), show that
    $I(\alpha) = \sqrt{\pi / \alpha}$.
    \item Differentiate $I(\alpha)$ with respect to $\alpha$ under the
    integral sign to derive $\int_{-\infty}^{\infty} x^2
    e^{-\alpha x^2}\,dx = \frac{1}{2}\sqrt{\pi/\alpha^3}$.
    \item Use parts (a) and (b) with $\alpha = 1/2$ to re-derive
    $\E[Z^2] = 1$ for $Z \sim \mathcal{N}(0,1)$.
\end{enumerate}
\textit{Hint:} Differentiation under the integral sign is justified by
dominated convergence, since $x^2 e^{-\alpha x^2}$ is integrable for
$\alpha > 0$.
\end{exercise}

\begin{exercise}[Bandit Reward Distributions]
\textnormal{(\(\ast\ast\ast\))} \\
A two-armed bandit has arm~1 with reward
$R_1 \sim \mathrm{Bernoulli}(0.6)$ and arm~2 with reward
$R_2 \sim \mathcal{N}(0.5, 0.25)$.
\begin{enumerate}[nosep,label=(\alph*)]
    \item Compute $\E[R_1]$, $\Var(R_1)$, $\E[R_2]$, and $\Var(R_2)$.
    \item Which arm has the higher expected reward?
    \item Suppose you pull arm~2 and observe $r_2 = 1.3$. Compare the
    density $f_{R_2}(1.3)$ with the probability $\Prob(R_1 = 1) = 0.6$.
    Explain why comparing a density value with a probability is
    invalid, even though both are numbers associated with the observation.
\end{enumerate}
\end{exercise}

%% ============================================================
%% SECTION 7: SUMMARY
%% ============================================================

\section{Summary}
\label{sec:summary}

\begin{itemize}
    \item A \emph{random variable} $X \colon \Omega \to \R$ assigns a
    numerical value to each outcome; its distribution is fully
    characterized by the CDF $F_X(x) = \Prob(X \leq x)$.
    \item A \emph{discrete} random variable takes countably many values;
    its distribution is described by the PMF
    $p_X(k) = \Prob(X = k)$.
    \item The \emph{Bernoulli}$(p)$ distribution models a single trial
    with success probability $p$; $\E[X] = p$, $\Var(X) = p(1-p)$.
    \item The \emph{Binomial}$(n, p)$ distribution counts successes in
    $n$ independent trials; the indicator decomposition
    $X = \sum X_i$ gives $\E[X] = np$ and $\Var(X) = np(1-p)$.
    \item The \emph{Geometric}$(p)$ distribution counts the trial of the
    first success; $\E[X] = 1/p$ and $\Var(X) = (1-p)/p^2$.
    \item The \emph{Poisson}$(\lambda)$ distribution models rare event
    counts; it arises as the limit of $\mathrm{Bin}(n, \lambda/n)$ as
    $n \to \infty$, and satisfies $\E[X] = \Var(X) = \lambda$.
    \item A \emph{continuous} random variable has
    $\Prob(X = x) = 0$ for every $x$; its distribution is described by
    a PDF $f_X$ with
    $\Prob(a \leq X \leq b) = \int_a^b f_X(x)\,dx$.
    \item The \emph{Exponential}$(\lambda)$ distribution is the unique
    continuous memoryless distribution; $\E[X] = 1/\lambda$ and
    $\Var(X) = 1/\lambda^2$.
    \item The \emph{Normal}$(\mu, \sigma^2)$ distribution has the
    bell-shaped density normalized by the Gaussian integral
    $\int e^{-x^2/2}\,dx = \sqrt{2\pi}$; $\E[X] = \mu$ and
    $\Var(X) = \sigma^2$.
\end{itemize}

%% ============================================================
%% SECTION 8: FURTHER READING
%% ============================================================

\section{Further Reading}
\label{sec:further_reading}

\begin{itemize}
    \item \textbf{Lesson 11} develops expectation, variance, covariance,
    and joint distributions in full generality, including the linearity
    of expectation and the tower property of conditional expectation.
    \item \textbf{Lesson 12} covers the law of large numbers and the
    central limit theorem, explaining why sample means converge and why
    the normal distribution arises so ubiquitously.
    \item \textbf{Phase 01, Lesson 4} (probability foundations) revisits
    random variables, expectation, and distributions in the
    measure-theoretic framework using Lebesgue integration.
    \item For a comprehensive catalog of over one hundred univariate
    distributions and their relationships, see
    \citet[Appendix A]{blitzstein2019}.
    \item The Poisson distribution was introduced by Sim\'{e}on Denis
    Poisson in 1837 to model wrongful convictions in court; the Gaussian
    distribution was discovered independently by de Moivre (1733), Gauss
    (1809), and Laplace (1812).
    \item \citet[Chapter 4]{feller1968} provides a deep treatment of
    generating functions for discrete distributions, including elegant
    proofs of the Poisson limit theorem and moment calculations.
\end{itemize}

\bibliography{probability-refs}

\end{document}
